\section{Discussion and open problems}\label{sec:discusion}
The comparison of two cut-offs gives the relative expansion
\(Q(n)/Q_{\rm pred}(n)=1+\cerr/n+o(1/n)\). Equation~\eqref{eq:main} determines \(\cerr\) from
the profile and the three boundary terms. Two distinct questions remain.

The first is to certify the sign and the digits of \(\cerr\). The global interval obtained contains
zero. Narrowing it requires controlling the full nodal values and the signed mean \(J\). One route
would be to prove a uniform remainder for the spatial quadrature of the sixteen physical
contributions. This remainder would have to control the diagonal exclusions and the kernels
containing the resolvent, in particular the contraction \(Q_2\). In the continuum approach, an
enclosure that preserves the dependence between signed contributions would have to improve on the
half-width observed at the tenth GL20 node.

The second concerns the remainder in the count. After subtracting \(\cerr/n\) from \(R_Q(n)\),
an error \(o(1/n)\) remains. Its order is unknown, as is whether the period-four oscillation
observed at small sizes persists in it. Pointwise control of arrival and stability under arbitrary
vetoes are stronger estimates than those used here.
